% =====================================================================
%  STAT 482 -- Bayesian Statistics, Spring 2026
%  Beamer Presentation Template for Paper Presentations
%
%  INSTRUCTIONS FOR STUDENTS
%  --------------------------
%  1. Fill in the four \newcommand lines in the CUSTOMISATION block.
%  2. Replace every [TODO: ...] placeholder with your own content.
%  3. Keep each slide to 5-6 bullet points maximum.
%  4. Do NOT remove [fragile] from the Stan code frame (Slide 8).
%  5. Compile with pdflatex twice for correct page numbers.
%  6. Recommended: 8-10 slides for a 15-minute talk.
%
%  Compilation: pdflatex (run twice)
% =====================================================================

\documentclass[aspectratio=169,11pt]{beamer}

\usetheme{Madrid}

\usepackage{amsmath,amssymb}
\usepackage{booktabs}
\usepackage{graphicx}
\usepackage{listings}

% ---- QU Colours ----------------------------------------------------------
\definecolor{qumaroon}{RGB}{134,38,51}
\definecolor{qugold}{RGB}{191,156,85}
\definecolor{qulightgold}{RGB}{250,244,225}
\definecolor{stanblue}{RGB}{30,100,180}

\setbeamercolor{palette primary}    {bg=qumaroon,           fg=white}
\setbeamercolor{palette secondary}  {bg=qumaroon!85!black,  fg=white}
\setbeamercolor{palette tertiary}   {bg=qugold,             fg=white}
\setbeamercolor{palette quaternary} {bg=qumaroon!90!black,  fg=white}
\setbeamercolor{structure}          {fg=qumaroon}
\setbeamercolor{frametitle}         {bg=qumaroon,           fg=white}
\setbeamercolor{title}              {bg=qumaroon,           fg=white}
\setbeamercolor{block title}        {bg=qumaroon!80!black,  fg=white}
\setbeamercolor{block body}         {bg=qumaroon!6}
\setbeamercolor{block title alerted}{bg=qugold!85!black,    fg=white}
\setbeamercolor{block body alerted} {bg=qulightgold}
\setbeamercolor{block title example}{bg=stanblue!75,        fg=white}
\setbeamercolor{block body example} {bg=stanblue!8}
\setbeamercolor{item}               {fg=qumaroon}
\setbeamercolor{subitem}            {fg=qugold!80!black}
\setbeamercolor{footline}           {bg=qumaroon!90!black,  fg=white}

% ---- Fonts ---------------------------------------------------------------
\setbeamerfont{title}      {size=\large,     series=\bfseries}
\setbeamerfont{frametitle} {size=\normalsize,series=\bfseries}
\setbeamerfont{block title}{size=\small,     series=\bfseries}

% ---- Itemize / enumerate style -------------------------------------------
\setbeamertemplate{itemize item}   {\color{qumaroon}$\blacktriangleright$}
\setbeamertemplate{itemize subitem}{\color{qugold}$\bullet$}
\setbeamertemplate{enumerate item}
  {\textcolor{qumaroon}{\bfseries\insertenumlabel.}}
\setbeamertemplate{navigation symbols}{}

% ---- Custom footer -------------------------------------------------------
\setbeamertemplate{footline}{%
  \leavevmode
  \begin{beamercolorbox}[wd=0.34\paperwidth,ht=2.2ex,dp=1ex,
                         leftskip=4pt]{footline}%
    \small STAT 482 -- Bayesian Statistics
  \end{beamercolorbox}%
  \begin{beamercolorbox}[wd=0.40\paperwidth,ht=2.2ex,dp=1ex,
                         center]{footline}%
    \small\insertshortauthor
  \end{beamercolorbox}%
  \begin{beamercolorbox}[wd=0.26\paperwidth,ht=2.2ex,dp=1ex,
                         rightskip=4pt,right]{footline}%
    \small\insertframenumber\,/\,\inserttotalframenumber
  \end{beamercolorbox}%
  \vskip0pt
}

% ---- R / Stan listing style ----------------------------------------------
\lstset{
  language=R,
  basicstyle=\ttfamily\tiny,
  keywordstyle=\color{stanblue}\bfseries,
  commentstyle=\color{green!50!black}\itshape,
  stringstyle=\color{red!60!black},
  breaklines=true,
  backgroundcolor=\color{gray!7},
  showstringspaces=false,
  tabsize=2,
  frame=none,
  morekeywords={stan,sampling,extract,library,rstan,ggmcmc,hdi,HDInterval}
}

% =====================================================================
%  STUDENT CUSTOMISATION -- EDIT THESE FOUR LINES
% =====================================================================
\newcommand{\papertitle}{[TODO: Full Paper Title]}
\newcommand{\paperref}{[TODO: Author(s), Journal, Year]}
\newcommand{\studentname}{[TODO: Your Full Name]}
\newcommand{\paperdomain}{[TODO: e.g.\ Clinical / Ecology / Economics]}
% =====================================================================

\title[\papertitle]{\papertitle}
\subtitle{STAT 482 -- Bayesian Statistics | Spring 2026}
\author[\studentname]{%
  \texorpdfstring{\studentname\vskip4pt\small\paperref}{\studentname}
}
\institute{Department of Mathematics and Statistics \\ Qatar University}
\date{Spring 2026}

% =====================================================================
\begin{document}
% =====================================================================

% ------------------------------------------------------------------
%  SLIDE 1 -- Title
% ------------------------------------------------------------------
\begin{frame}
  \titlepage
\end{frame}

% ------------------------------------------------------------------
%  SLIDE 2 -- Roadmap
% ------------------------------------------------------------------
\begin{frame}{Presentation Roadmap}

  \begin{columns}[T]
    \begin{column}{0.52\textwidth}
      \begin{alertblock}{Paper in One Sentence}
        \small
        [TODO: Summarise the paper's contribution in one sentence
        for a non-specialist audience.]
      \end{alertblock}
      \vspace{6pt}
      \textbf{This presentation covers:}
      \begin{enumerate}\small
        \item Problem and motivation
        \item The Bayesian model
        \item Computational method
        \item Key findings
        \item My R/RStan replication
        \item Lessons learned
      \end{enumerate}
    \end{column}
    \begin{column}{0.44\textwidth}
      \begin{block}{Paper at a Glance}
        \small
        \textbf{Authors:} [TODO]\\
        \textbf{Journal:} [TODO]\\
        \textbf{Year:} [TODO]\\
        \textbf{Method:} [TODO: e.g.\ RStan + HMC]\\
        \textbf{Domain:} \paperdomain\\
        \textbf{Key tool:} R / RStan
      \end{block}
    \end{column}
  \end{columns}

\end{frame}

% ------------------------------------------------------------------
%  SLIDE 3 -- Problem and Motivation
% ------------------------------------------------------------------
\begin{frame}{Step 1 -- Problem and Motivation}

  \begin{block}{The Research Question}
    \small
    [TODO: State the research question in one or two sentences.
    What problem does the paper solve? Who cares about the answer?]
  \end{block}

  \vspace{5pt}

  \begin{columns}[T]
    \begin{column}{0.48\textwidth}
      \textbf{Data and Setting:}
      \begin{itemize}\small
        \item [TODO: dataset description]
        \item [TODO: sample size, outcomes measured]
        \item [TODO: study design -- RCT, observational, etc.]
      \end{itemize}
    \end{column}
    \begin{column}{0.48\textwidth}
      \textbf{Why Bayesian?}
      \begin{itemize}\small
        \item [TODO: reason 1]
        \item [TODO: reason 2]
        \item [TODO: reason 3]
      \end{itemize}
    \end{column}
  \end{columns}

  \vspace{5pt}

  \begin{alertblock}{Where classical methods fall short}
    \small
    [TODO: 1--2 sentences on why a frequentist approach is
    insufficient or less informative for this specific problem.]
  \end{alertblock}

\end{frame}

% ------------------------------------------------------------------
%  SLIDE 4 -- The Bayesian Model
% ------------------------------------------------------------------
\begin{frame}{Step 2 -- The Bayesian Model}

  \begin{block}{Likelihood}
    \[
      [TODO:\quad p(\mathbf{y}\mid\boldsymbol{\theta}) = \ldots]
    \]
    \small
    [TODO: What are $\mathbf{y}$ (observed data) and
    $\boldsymbol{\theta}$ (parameters) in this paper's context?]
  \end{block}

  \vspace{3pt}

  \begin{block}{Prior Distribution(s)}
    \[
      [TODO:\quad\theta_1 \sim \mathcal{N}(\ldots,\ldots)]
      \qquad
      [TODO:\quad\theta_2 \sim \ldots]
    \]
    \small
    \textit{Justification:}
    [TODO: one sentence on why these priors are appropriate.]
  \end{block}

  \vspace{3pt}

  \begin{exampleblock}{Posterior -- Bayes' Theorem}
    \[
      \pi(\boldsymbol{\theta}\mid\mathbf{y})
      \;\propto\;
      \underbrace{p(\mathbf{y}\mid\boldsymbol{\theta})}_{\text{likelihood}}
      \;\times\;
      \underbrace{\pi(\boldsymbol{\theta})}_{\text{prior}}
    \]
    \small
    [TODO: What does this posterior represent in context?]
  \end{exampleblock}

\end{frame}

% ------------------------------------------------------------------
%  SLIDE 5 -- Model Structure (adapt to your paper)
% ------------------------------------------------------------------
\begin{frame}{Step 2 (continued) -- Model Structure}

  %  Adapt this slide to your assigned paper:
  %  - Hierarchical model  -> keep the three-level structure below.
  %  - Regression model    -> show the prior sensitivity panel instead.
  %  - Methods paper (HMC) -> replace with Hamiltonian dynamics sketch.

  \begin{columns}[T]
    \begin{column}{0.50\textwidth}
      \textbf{[TODO: Name of model component / level]}
      \vspace{4pt}
      \begin{exampleblock}{\small Level 1 -- [TODO: label]}
        \small
        $[TODO:\;\hat{\theta}_k \mid \theta_k
          \sim \mathcal{N}(\theta_k,\,\sigma_k^2)]$
      \end{exampleblock}
      \vspace{2pt}
      \begin{exampleblock}{\small Level 2 -- [TODO: label]}
        \small
        $[TODO:\;\theta_k \mid \mu,\tau
          \sim \mathcal{N}(\mu,\,\tau^2)]$
      \end{exampleblock}
      \vspace{2pt}
      \begin{exampleblock}{\small Level 3 -- [TODO: label]}
        \small
        $[TODO:\;\mu \sim \ldots,\quad \tau \sim \ldots]$
      \end{exampleblock}
    \end{column}
    \begin{column}{0.46\textwidth}
      \begin{block}{Key Insight}
        \small
        [TODO: What does this model structure achieve that a simpler
        model could not? e.g.\ partial pooling, prior sensitivity,
        HMC efficiency.]
      \end{block}
      \vspace{5pt}
      \textbf{Parameter interpretation:}
      \begin{itemize}\small
        \item $\mu$ = [TODO]
        \item $\tau$ = [TODO]
        \item [TODO: other parameters]
      \end{itemize}
    \end{column}
  \end{columns}

\end{frame}

% ------------------------------------------------------------------
%  SLIDE 6 -- Computation and Diagnostics
% ------------------------------------------------------------------
\begin{frame}{Step 3 -- Computation and Diagnostics}

  \begin{columns}[T]
    \begin{column}{0.50\textwidth}
      \textbf{RStan setup:}
      \begin{itemize}\small
        \item Algorithm: [TODO: HMC/NUTS or other]
        \item Chains: [TODO] \quad Iterations: [TODO]
        \item Warm-up: [TODO]
        \item \texttt{adapt\_delta}: [TODO]
        \item Parameterisation: [TODO: centred / non-centred]
      \end{itemize}
      \vspace{5pt}
      \textbf{Convergence checks:}
      \begin{itemize}\small
        \item $\hat{R}$: [TODO: value; threshold $< 1.01$]
        \item ESS: [TODO: adequate?]
        \item Divergent transitions: [TODO: any?]
        \item Trace plots: [TODO: mixed well?]
      \end{itemize}
    \end{column}
    \begin{column}{0.46\textwidth}
      \begin{block}{Trace Plot -- [TODO: parameter name]}
        \vspace{28pt}
        \centering\small
        \textit{[Insert your trace plot here]}\\
        \tiny
        Good mixing: chains overlap and\\
        show no long drifts or stuck regions.
        \vspace{28pt}
      \end{block}
    \end{column}
  \end{columns}

\end{frame}

% ------------------------------------------------------------------
%  SLIDE 7 -- Key Findings
% ------------------------------------------------------------------
\begin{frame}{Step 4 -- Key Findings}

  \begin{alertblock}{Main Result}
    \small
    [TODO: State the most important finding with a number.
    Example: ``The posterior mean treatment effect is $-14.2$ min
    (95\% CrI: $-22.1$, $-6.3$), indicating a likely reduction
    in PACU discharge time.'']
  \end{alertblock}

  \vspace{5pt}

  \begin{columns}[T]
    \begin{column}{0.50\textwidth}
      \textbf{Posterior summaries:}
      \vspace{3pt}
      {\small
      \begin{tabular}{lcc}
        \toprule
        Parameter & Post.\ mean & 95\% CrI \tabularnewline
        \midrule
        TODO: name & val & (lo,\ hi) \tabularnewline
        TODO: name & val & (lo,\ hi) \tabularnewline
        TODO: name & val & (lo,\ hi) \tabularnewline
        \bottomrule
      \end{tabular}}
    \end{column}
    \begin{column}{0.46\textwidth}
      \begin{exampleblock}{Posterior Probability Statement}
        \small
        $P(\text{[TODO: effect > threshold]} \mid \mathbf{y}) = $ [TODO]\par
        \textit{In plain language:}\par
        [TODO: interpret this for a non-statistician]
      \end{exampleblock}
      \vspace{4pt}
      \small
      [TODO: One further finding -- robustness, heterogeneity,
      or comparison to a frequentist estimate.]
    \end{column}
  \end{columns}

\end{frame}

% ------------------------------------------------------------------
%  SLIDE 8 -- Replication   [fragile is REQUIRED for lstlisting]
% ------------------------------------------------------------------
\begin{frame}[fragile]{Step 5 -- My R / RStan Replication}

  \begin{columns}[T]
    \begin{column}{0.50\textwidth}
      \textbf{What I reproduced:}
      \begin{itemize}\small
        \item [TODO: model name]
        \item [TODO: real or simulated data?]
        \item [TODO: key result reproduced]
      \end{itemize}

      \vspace{4pt}
      \textbf{Stan \texttt{model\{\}} block:}
\begin{lstlisting}
// [TODO: paste your Stan model block]
// Show the parameters{} and model{} sections
parameters {
  real beta1;     // treatment effect
  real<lower=0> sigma;
}
model {
  beta1 ~ normal(0, 10);  // prior
  y ~ normal(beta0 + beta1 * x, sigma);
}
\end{lstlisting}
    \end{column}

    \begin{column}{0.46\textwidth}
      \begin{block}{My Key Output}
        \vspace{28pt}
        \centering\small
        \textit{[Insert your main plot or table here]}\\
        \tiny
        e.g.\ posterior density, forest plot,\\
        credible interval plot, shrinkage plot
        \vspace{28pt}
      \end{block}
      \vspace{4pt}
      \small
      \textbf{Matches the paper?}\par
      [TODO: Qualitative comparison -- direction,
      magnitude, and uncertainty.]
    \end{column}
  \end{columns}

\end{frame}

% ------------------------------------------------------------------
%  SLIDE 9 -- Lessons Learned
% ------------------------------------------------------------------
\begin{frame}{Step 6 -- Lessons Learned}

  \begin{columns}[T]
    \begin{column}{0.58\textwidth}
      \textbf{Three things I understand better now:}
      \vspace{5pt}
      \begin{enumerate}\small
        \item \textbf{[TODO: Statistical concept]}\\
              \small
              [TODO: What you understand now that you did not
              before reading the paper -- be specific.]
        \vspace{4pt}
        \item \textbf{[TODO: Link to STAT 482]}\\
              \small
              [TODO: Connect to a lecture topic -- prior choice,
              MCMC diagnostics, hierarchical models, \ldots]
        \vspace{4pt}
        \item \textbf{[TODO: Practical takeaway]}\\
              \small
              [TODO: One thing you will do differently
              in your own Bayesian analyses.]
      \end{enumerate}
    \end{column}
    \begin{column}{0.38\textwidth}
      \begin{alertblock}{Open Question}
        \small
        [TODO: One aspect the paper does not address that
        you think would strengthen or extend the analysis.]
      \end{alertblock}
      \vspace{6pt}
      \begin{exampleblock}{Connection to This Course}
        \small
        [TODO: Name the STAT 482 lecture or topic
        this paper most directly illustrates or extends.]
      \end{exampleblock}
    \end{column}
  \end{columns}

\end{frame}

% ------------------------------------------------------------------
%  SLIDE 10 -- References and Q&A
% ------------------------------------------------------------------
\begin{frame}{References and Thank You}

  \begin{block}{Assigned Paper}
    \small
    \paperref\\
    \vspace{3pt}
    DOI: \texttt{[TODO: insert DOI]}
  \end{block}

  \vspace{5pt}

  \begin{block}{R Packages Used in My Replication}
    \small
    \begin{itemize}
      \item \texttt{rstan} -- Bayesian model fitting via Stan
      \item \texttt{ggmcmc} -- MCMC diagnostic plots
      \item \texttt{HDInterval} -- Highest density credible intervals
      \item [TODO: any additional packages used]
    \end{itemize}
  \end{block}

  \vspace{8pt}

  \begin{center}
    {\large\bfseries\color{qumaroon} Thank you -- Questions Welcome}\\
    \small\textit{R/Stan code available upon request.}
  \end{center}

\end{frame}

% =====================================================================
\end{document}
